\documentclass{article}
\usepackage[T1]{fontenc}
\usepackage{lmodern}
\usepackage{graphicx}
\usepackage{amsmath,amssymb}
\usepackage[a4paper,margin=2cm]{geometry}
\usepackage[absolute,overlay]{textpos}
\setlength{\TPHorizModule}{1mm}
\setlength{\TPVertModule}{1mm}
\usepackage[indonesian]{babel}
\usepackage{hyperref}
\setlength{\emergencystretch}{2em}
% unit: Halaman 6

\begin{document}

% === Halaman 6 (python/native) ===

• Sayangnya, algoritma knapsack sederhana di atas hanya dapat digunakan 

untuk enkripsi, tetapi tidak untuk dekripsi. 

• Misalnya, jika diberikan kriptogram = 32, maka tentukan b1, b2, …, b6 

sedemikian sehingga 


32= b1 + 5b2 + 6b3 + 11b4 + 14b5 + 20b6 


(2)

• Solusi persamaan (2) ini tidak dapat dipecahkan dalam orde waktu polinomial 

dengan semakin besarnya n (dengan catatan barisan bobot tidak dalam urutan 

menaik). 

• Namun, hal inilah yang dijadikan sebagai kekuatan  algoritma knapsack.

6

\end{document}
